\section{Methodology}
\subsection{Numerical method}
For arbitrary viscosity profiles, \autoref{eq:EkmanEquation} generally does not admit a closed-form analytical solution. The equation was therefore solved numerically as a boundary value problem using the \texttt{solve\_bvp} routine available in \texttt{SciPy} \citep{virtanen2020}. 

For the SUBC model, it was shown that the equation could be rewritten for a dimensionless layer thickness $\varphi$. Thus, using an arbitrary height-normalized height-dependent eddy viscosity profile, $\tilde\nu(z)$, one can write a simpler differential equation. Separating the equation into the zonal and meridional velocity components, one obtains 
\begin{subequations}
    \begin{align}
    \frac{\mathrm{d}}{\mathrm{d} z}
    \left(
    \tilde\nu(z)\frac{\mathrm{d} u}{\mathrm{d} z}
    \right)
    &=
    -\varphi^2 v,
    \\
    \frac{\mathrm{d}}{\mathrm{d} z}
    \left(
    \tilde\nu(z)\frac{\mathrm{d} v}{\mathrm{d} z}
    \right)
    &=
    \varphi^2 (u-u_g),
    \label{eq:Method}
\end{align}
\end{subequations}
where the wind aloft was set as a purely zonal wind $u_g$. To apply the numerical boundary value solver, the equations were rewritten as a first-order system by introducing the new variables
\begin{equation}
    u^\prime = \frac{\mathrm{d} u}{\mathrm{d} z},
    \qquad
    v^\prime = \frac{\mathrm{d} v}{\mathrm{d} z}.
\end{equation}
Thus one can expand \autoref{eq:Method} into four coupled first-order equations
\begin{align}
    \begin{cases} 
    \frac{\mathrm{d} u}{\mathrm{d} z}
    &= u^\prime,
    \\
    \frac{\mathrm{d} u^\prime}{\mathrm{d} z}
    &=
    \frac{-\varphi^2 v - \tilde\nu^\prime u^\prime}
    {\tilde\nu},
    \\
    \frac{\mathrm{d} v}{\mathrm{d} z}
    &= v^\prime,
    \\
    \frac{\mathrm{d} v^\prime}{\mathrm{d} z}
    &=
    \frac{\varphi^2(u-u_g)-\tilde\nu^\prime v^\prime}
    {\tilde\nu},
    \end{cases}
    \label{eq:modelsolver}
\end{align}
with the boundary conditions
\begin{subequations}
    \begin{align}
    u(\varphi=0) &= 0,
    &
    u^\prime(\varphi=1) &= 0,
    \\
    v(\varphi=0) &= 0,
    &
    v^\prime(\varphi=1) &= 0.
\end{align}
\end{subequations}
If the viscosity profile, $\tilde\nu$, obtains a value of zero in the domain, then the coupled \autoref{eq:modelsolver} becomes unsolvable since we divide by $\hat{\nu}=0$. In this scenario one has to introduce a constant background viscosity, $\epsilon$ as
\begin{equation}
    \hat{\nu}\rightarrow\hat{\nu}+\epsilon,
\end{equation}
where $\hat{\nu}/\epsilon\gg1$ was assumed. This background viscosity was thus applied for the linear increasing model and the parabolic model since both have zero viscosity at the surface.

The SUBC model equations were solved iteratively on a vertical grid. The solver required an initial guess, which was chosen based on the initial setup of the system. The initial guess was therefore set as the classical solution for a constant viscosity. The final solution was found when the iterative solution reached a steady state for the differential equations with the imposed boundary conditions.